\section{Chapter 11}

\begin{verse}
Thirty spokes converge in the hub\\
But the essence of the wheel is the emptiness of the hub.\\
The vessels are made of clay\\
But it is the inner emptiness that makes the essence of the vessel\\
The wall with windows and doors form a house\\
But it is their emptiness that makes its essence.\\
In general: the being serves as the useful means\\
The essence lies in non-being in emptiness.
\end{verse}

Various images are used to express the idea that the essence of the material and the sensible lies in the immaterial and the invisible, that the ``fullness'' is ordered to ``emptiness'': a reference to the metaphysical plane (``emptiness''=transcendence) and a reference concordant with the nature of non-acting. The wheel with thirty spokes was that of an ancient sacred chariot.

In the esoteric interpretation typical to operative Taoism this chapter was also put in relation with the distillation of the subtle from the dense, the second being ordered to the first and having to be resolved in the first (in the yang soul) in the ``solution of the form''.

For the ambivalence of the Tao, see the expression: ``It is as heavy as a stone, as light as a feather.''

\paragraph{Chinese text and literal translation}

\textbf{Chapter 11 (\begin{CJK*}{UTF8}{bsmi}第十一章\end{CJK*})}

\columnratio{.4}
\begin{paracol}{2}
\begin{verse}
\begin{CJK*}{UTF8}{bsmi}
三十輻共一轂，\\
當其無，有車之用。\\
埏埴以為器，\\
當其無，有器之用。\\
鑿戶牖以為室，\\
當其無，有室之用。\\
故有之以為利，\\
無之以為用。
\end{CJK*}
\end{verse}
\switchcolumn
\begin{verse}
Put thirty spokes together to one hub,\\
The original empty space makes the use of wheel.\\
Knead clay into vessels,\\
The original empty space makes the use of vessel.\\
Shape door and windows for a house/room,\\
The original empty space makes the use of house/room.\\
So the things that are made are only conditions,\\
What we are using is still the original empty space.
\end{verse}
\end{paracol}


\flrightit{Posted on 2020-10-11 by Lao Tzu}

